% sections/04-experiments.tex

\section{Experiments}
\label{sec:experiments}

All experiments use synthetic case count data generated on the real administrative
structure of Vietnam (63 provinces), with standard surveillance parameters:
spatial window up to 10\% of $\Ntot$, temporal window 14 days,
$|\Wcal| \approx 6{,}000$ candidate windows, $M = 999$ Monte Carlo simulations.
No real disease data is used; province boundaries and adjacency are real.
All code is in Python, CPU-only, running on \placeholder{hardware spec}.

% ── Experiment 1 ──────────────────────────────────────────────────────────────

\subsection{Experiment 1: Sensitivity Surface}
\label{sec:exp1}

\paragraph{Goal.}
Visualize $\sens$ as a function of $\kw$ to empirically validate the functional
form of Theorem~\ref{thm:sensitivity} and identify the tightness gap between the
upper bound and the empirical maximum sensitivity.

\paragraph{Setup.}
\placeholder{For each $\kw \in \{2, 3, 5, 10, 20, 50, 100, 200\}$, generate 10{,}000
random neighboring dataset pairs. Compute $|\LLR(w, D) - \LLR(w, D')|$ for each pair.
Report the empirical maximum alongside the bound $f(\kw)$.}

\paragraph{Results.}
\placeholder{Figure~\ref{fig:sensitivity-surface} shows the sensitivity surface.
The bound $f(\kw)$ \todo{is/is not} tight within a factor of $X$ across all densities
tested. The largest gap occurs at $\kw = \todo{?}$.}

\begin{figure}[ht]
    \centering
    \placeholder{Figure: sensitivity surface $\Delta(\kw)$ vs. $\kw$.
    Two curves: empirical maximum (dots) and bound $f(\kw)$ (line).
    Log scale on y-axis. Key reference points: Ha Giang ($\kw=3$), Hanoi ($\kw=45$),
    Dong Nai ($\kw=200$).
    File: figures/exp1-sensitivity-surface.pdf}
    \caption{Sensitivity bound $\sens$ as a function of window density $\kw$.
             The upper bound from Theorem~\ref{thm:sensitivity} (solid) is compared to
             the empirical maximum sensitivity (dots) over 10{,}000 random dataset pairs.}
    \label{fig:sensitivity-surface}
\end{figure}

% ── Experiment 2 ──────────────────────────────────────────────────────────────

\subsection{Experiment 2: Detection Quality Under Privacy}
\label{sec:exp2}

\paragraph{Goal.}
Compare detection quality of the proposed exponential mechanism (Theorem~\ref{thm:mechanism})
against the Laplace baseline across a range of privacy budgets $\eps \in \{0.1, 0.5, 1.0, 2.0, 5.0\}$.

\paragraph{Setup.}
\placeholder{Inject a synthetic outbreak in a dense province ($\kw = 200$) and a sparse
province ($\kw = 3$). For each mechanism and each $\eps$, run 1{,}000 trials.
Measure: (1) fraction of trials where $\wstar$ matches the injected cluster,
(2) expected $\LLR(\wstar)$ relative to the non-private optimum.}

\paragraph{Results.}
\placeholder{Table~\ref{tab:detection} shows detection rates.
The exponential mechanism maintains $>\!90\%$ detection rate for the dense cluster
at $\eps = \todo{?}$, where the Laplace baseline drops to $\todo{?}\%$.
For the sparse cluster, both mechanisms perform similarly because the sensitivity
is high regardless.}

\begin{table}[ht]
    \centering
    \caption{Detection rate (\% of trials where $\wstar$ = injected cluster) for the
             exponential mechanism (Theorem~\ref{thm:mechanism}) vs. the Laplace
             baseline, at varying $\eps$.}
    \label{tab:detection}
    \begin{tabular}{@{}lcccccc@{}}
        \toprule
        & \multicolumn{5}{c}{Privacy budget $\eps$} \\
        \cmidrule(l){2-6}
        Mechanism & 0.1 & 0.5 & 1.0 & 2.0 & 5.0 \\
        \midrule
        Laplace (dense, $\kw=200$)      & \placeholder{?} & \placeholder{?} & \placeholder{?} & \placeholder{?} & \placeholder{?} \\
        Exponential (dense, $\kw=200$)  & \placeholder{?} & \placeholder{?} & \placeholder{?} & \placeholder{?} & \placeholder{?} \\
        \midrule
        Laplace (sparse, $\kw=3$)       & \placeholder{?} & \placeholder{?} & \placeholder{?} & \placeholder{?} & \placeholder{?} \\
        Exponential (sparse, $\kw=3$)   & \placeholder{?} & \placeholder{?} & \placeholder{?} & \placeholder{?} & \placeholder{?} \\
        \bottomrule
    \end{tabular}
\end{table}

% ── Experiment 3 ──────────────────────────────────────────────────────────────

\subsection{Experiment 3: Surveillance Horizon Under Composition}
\label{sec:exp3}

\paragraph{Goal.}
Validate the $\Tmax$ formula from Theorem~\ref{thm:composition} empirically and
compare the operational horizon of zCDP composition against naive composition across
a range of $\rhow$ values.

\paragraph{Setup.}
\placeholder{For each $\rhow \in \{0.001, 0.005, 0.01, 0.05, 0.1\}$ and target budget
$\eps_{\mathrm{target}} \in \{0.5, 1.0, 2.0\}$, compute $\Tmax$ from the formula and
compare to the exact quadratic solution.
Report the approximation error as a function of $\rhow$.}

\paragraph{Results.}
\placeholder{Figure~\ref{fig:horizon} shows $\Tmax$ as a function of $\rhow$.
The approximation is within $\todo{?}\%$ of the exact solution for $\rhow < 0.05$.
At $\rhow = 0.01$, $\eps = 1$, $\delta = 10^{-6}$: $\Tmax \approx 180$ weeks vs.
10 weeks under naive composition.}

\begin{figure}[ht]
    \centering
    \placeholder{Figure: $\Tmax$ vs. $\rhow$ for three target budgets.
    Two curves per budget: zCDP formula (solid) and naive composition (dashed).
    Log scale on y-axis (weeks).
    File: figures/exp3-surveillance-horizon.pdf}
    \caption{Maximum surveillance duration $\Tmax$ (weeks) as a function of the
             per-week zCDP cost $\rhow$, for target budgets
             $\eps \in \{0.5, 1.0, 2.0\}$.
             Solid: Theorem~\ref{thm:composition} (zCDP).
             Dashed: naive composition.
             $\delta = 10^{-6}$ throughout.}
    \label{fig:horizon}
\end{figure}
